\documentclass[11pt,reqno]{amsart}

\usepackage[T1]{fontenc}
\usepackage[utf8]{inputenc}
\usepackage{lmodern}
\usepackage{amsmath,amssymb,amsfonts,mathtools,mathrsfs}
\usepackage{amsthm}
\usepackage{microtype}
\usepackage{enumitem}
\usepackage{booktabs}
\usepackage{array}
\usepackage{geometry}
\usepackage{hyperref}
\usepackage{bookmark}
\geometry{margin=1.08in}
\hypersetup{colorlinks=true,linkcolor=blue,citecolor=blue,urlcolor=blue}

\newtheorem{theorem}{Theorem}[section]
\newtheorem{proposition}[theorem]{Proposition}
\newtheorem{lemma}[theorem]{Lemma}
\newtheorem{corollary}[theorem]{Corollary}
\theoremstyle{definition}
\newtheorem{definition}[theorem]{Definition}
\newtheorem{hypothesis}[theorem]{Hypothesis}
\newtheorem{objective}[theorem]{Objective}
\theoremstyle{remark}
\newtheorem{remark}[theorem]{Remark}

\newcommand{\R}{\mathbb R}
\newcommand{\C}{\mathbb C}
\newcommand{\N}{\mathbb N}
\newcommand{\eps}{\varepsilon}
\newcommand{\li}{\operatorname{li}}
\newcommand{\Log}{\operatorname{Log}}
\newcommand{\Ei}{\operatorname{Ei}}
\newcommand{\Res}{\operatorname*{Res}}
\newcommand{\Mell}{*_{\rm Mel}}
\newcommand{\RVM}{\mathrm{RVM}}
\newcommand{\crit}{\mathrm{crit}}
\newcommand{\linez}{\mathrm{line}}
\newcommand{\pts}{\mathrm{pts}}
\newcommand{\pathc}{\mathrm{path}}
\newcommand{\dd}{\,d}

\title[The midpoint RVM synthetic zeta system]{The midpoint Riemann--von Mangoldt synthetic zeta system:\\
canonical products, exponential sums, prime counts, Mertens transforms, and ordinate transport}
\author{V\'ictor Casa\~na-Giner}
\address{Bio Logic Crop Science S.L.}
\thanks{ORCID: 0000-0002-8252-2608.}
\date{4 October 2026}

\begin{document}
\begin{abstract}
We isolate a deterministic zeta-like system generated by the midpoint inverse of the smooth Riemann--von Mangoldt counting law.  Its synthetic nontrivial zeros are prescribed to be
$\rho_k^{\RVM}=\tfrac12+i\gamma_k^{\RVM}$; the critical real part is therefore an input of the construction and carries no implication for the genuine zeros of the Riemann zeta function.  We prove normal convergence of the associated canonical product, its functional symmetry, the exact synthetic zero divisor, a uniform logarithmic bound for the exponential sums of the midpoint ordinates, compact-uniform existence of the corresponding explicit-formula prime-power count, and exact finite M\"obius inversion to a synthetic prime count.  We separate these theorems from the analytic hypotheses required for a square-root-scale prime-counting estimate and for a continuous-to-atomic realization on the integers.

On the Mertens side we define the transform $\widehat M_{\RVM}^{\crit}=1/\zeta_{\RVM}$, identify its pole abscissa as $1/2$ by construction, and give the exact quotient relating the synthetic transform to $1/\zeta$.  We then prove the midpoint transition formula by an Abel-regularized evaluation of the clock product.  It yields a constant vertical limit for $1/\zeta_{\RVM}$ on every fixed line to the right of $1/2$, proves that neither $\zeta_{\RVM}$ nor its reciprocal is an absolutely convergent ordinary Dirichlet series in any right half-plane, and justifies the symmetric unsmoothed Perron inverse with the bound $M_{\RVM}^{\crit}(x)=L+O_\varepsilon(x^{1/2+\varepsilon})$ for $x>1$.  Finally, we incorporate a continuous ordinate-transport argument: the midpoint product is deformed to the product supported on the genuine ordinates projected onto the critical line.  The entire product family converges jointly, its infinitesimal generator and total multiplier are explicit, and the endpoint identity with the classical zeta function is proved equivalent to the Riemann hypothesis.  The coefficient convolution proposed in the three-step note is therefore recorded as unavailable in the category of absolutely convergent ordinary Dirichlet series, rather than as an unproved route.
\end{abstract}
\maketitle

\tableofcontents

\section{Purpose, separation of objects, and status}

The object of this paper is a synthetic arithmetic system determined entirely by the midpoint smooth Riemann--von Mangoldt clock.  It is logically distinct from collision-energy, Catalan, and excursion programs.  None of those programs is used as a premise below.

Let the nontrivial zeros of the genuine Riemann zeta function in the upper half-plane, counted with multiplicity, be
\begin{equation}\label{eq:true-zeros}
 \rho_k=\beta_k+i\gamma_k,
 \qquad 0<\beta_k<1,
 \qquad 0<\gamma_1\le\gamma_2\le\cdots .
\end{equation}
No critical-line assumption is made in \eqref{eq:true-zeros}.  The synthetic midpoint ordinates are denoted
\begin{equation}\label{eq:synthetic-zero-notation}
 \rho_k^{\RVM}:=\frac12+i\gamma_k^{\RVM}.
\end{equation}
We never use $\tau_k$ as a substitute for $\gamma_k^{\RVM}$.

Three objects must be kept distinct throughout:
\begin{enumerate}[label=\textup{(\roman*)}]
\item $M_H(x)$ and $M_H^{\RVM}(x)$ are finite-height averages of phases of zero ordinates;
\item $M_{\RVM}^{\crit}$ is the inverse-Mellin object associated with $1/\zeta_{\RVM}$;
\item $M(x)=\sum_{n\le x}\mu(n)$ is the genuine Mertens function.
\end{enumerate}
The first object is not a Mertens summatory function.

\begin{center}
\begin{tabular}{p{0.34\textwidth}p{0.56\textwidth}}
\toprule
Statement & Status in this paper\\
\midrule
Midpoint clock and Lambert inversion & theorem / definition\\
Canonical product $P_{\RVM}^{\crit}$ and $\zeta_{\RVM}$ & theorem\\
$S_{\RVM}(N;u)\ll_{a,b}1+\log N$ & theorem\\
Compact-uniform existence of $J_{\RVM}$ & theorem\\
Finite inversion $J_{\RVM}^{*}\leftrightarrow\pi_{\RVM}$ & theorem\\
Square-root-scale estimate for $J_{\RVM},\pi_{\RVM}$ & hypothesis\\
Integer-weighted Euler realization & conditional on unit-cell and Mellin--Stieltjes control\\
Gaussian square-root bound for $M_{\RVM}^{\crit}$ & theorem conditional on stated vertical growth\\
Midpoint transition formula $Z_{\RVM}(1+w)$ and $Z_{\RVM}(1+\chi)$ & theorem\\
Symmetric unsmoothed Perron inverse $M_{\RVM}^{\crit}$ & theorem; $L+O_\varepsilon(x^{1/2+\varepsilon})$\\
Ordinary absolutely convergent Dirichlet expansion of $\zeta_{\RVM}^{\pm1}$ & impossible\\
Continuous ordinate transport & theorem\\
Endpoint identity $\zeta_{\linez}=\zeta$ & equivalent to RH\\
Coefficient identity with $\mu$ & endpoint criterion only; the RVM-factor convolution is unavailable\\
\bottomrule
\end{tabular}
\end{center}

\section{The midpoint Riemann--von Mangoldt clock}

Define
\begin{equation}\label{eq:clock}
 \mathcal N_{\RVM}(T)
 :=\frac{T}{2\pi}\log\frac{T}{2\pi e}+\frac78.
\end{equation}
The midpoint synthetic ordinates are the positive solutions of
\begin{equation}\label{eq:midpoint-condition}
 \mathcal N_{\RVM}(\gamma_k^{\RVM})=k-\frac12,
 \qquad k\ge1.
\end{equation}
The branch $W_0$ below denotes the principal branch of the Lambert $W$-function.

\begin{proposition}[Lambert inversion]\label{prop:lambert}
For every $k\ge1$,
\begin{equation}\label{eq:lambert}
 \boxed{
 \gamma_k^{\RVM}
 =\frac{2\pi(k-11/8)}{W_0((k-11/8)/e)}.}
\end{equation}
The quotient in \eqref{eq:lambert} is positive also for $k=1$.
\end{proposition}

\begin{proof}
Put $y=\gamma/(2\pi)$.  Equation \eqref{eq:midpoint-condition} is
\[
 y(\log y-1)=k-\frac{11}{8}.
\]
Writing $v=\log y-1$ gives $y=e^{v+1}$ and hence
$v e^v=(k-11/8)/e$.  Thus $v=W_0((k-11/8)/e)$ and
$y=(k-11/8)/W_0((k-11/8)/e)$.  For $k=1$ both numerator and the principal-branch value of $W_0$ are negative, so the quotient remains positive.
\end{proof}

The clock derivative is
\begin{equation}\label{eq:clock-derivative}
 \mathcal N_{\RVM}'(T)=\frac1{2\pi}\log\frac{T}{2\pi}.
\end{equation}
In particular
\begin{equation}\label{eq:clock-growth}
 \gamma_k^{\RVM}\asymp \frac{k}{\log k}.
\end{equation}

\begin{lemma}[Bounded displacement from the genuine ordinates]\label{lem:bounded-displacement}
With the enumeration \eqref{eq:true-zeros},
\begin{equation}\label{eq:bounded-displacement}
 \gamma_k-\gamma_k^{\RVM}=O(1).
\end{equation}
This is unconditional and includes multiplicities.
\end{lemma}

\begin{proof}
The Riemann--von Mangoldt formula gives
\[
 N_\zeta(T)=\mathcal N_{\RVM}(T)+O(\log(2+T)).
\]
At a zero ordinate $\gamma_k$, either endpoint convention for $N_\zeta$ gives
$N_\zeta(\gamma_k)=k+O(m_k)$, where $m_k$ is the multiplicity at that ordinate.  The same counting law gives $m_k=O(\log(2+\gamma_k))$.  Hence
\[
 \mathcal N_{\RVM}(\gamma_k)=k+O(\log(2+\gamma_k)).
\]
Together with \eqref{eq:midpoint-condition}, this gives a clock-coordinate discrepancy $O(\log\gamma_k)$.  By the mean-value theorem and \eqref{eq:clock-derivative}, whose size is $\asymp\log T$ at height $T$, the corresponding ordinate discrepancy is $O(1)$.  Finitely many initial ordinates are absorbed into the constant.
\end{proof}

\section{The pure critical RVM zeta function}

Put
\begin{equation}\label{eq:Pcrit}
 P_{\RVM}^{\crit}(s)
 :=\prod_{k\ge1}
 \left(1+\frac{s(s-1)}{1/4+(\gamma_k^{\RVM})^2}\right),
\end{equation}
\begin{equation}\label{eq:Gdef}
 G(s):=s(s-1)\Gamma(s/2)\pi^{-s/2},
\end{equation}
and
\begin{equation}\label{eq:zetaRVM}
 \xi_{\RVM}(s):=\frac12P_{\RVM}^{\crit}(s),
 \qquad
 \zeta_{\RVM}(s):=\frac{P_{\RVM}^{\crit}(s)}{G(s)}.
\end{equation}

\begin{theorem}[Canonical characterization of $\zeta_{\RVM}$]\label{thm:canonical}
The product \eqref{eq:Pcrit} converges absolutely and locally uniformly in $s\in\C$.  It defines an entire function satisfying
\begin{equation}\label{eq:P-symmetry}
 P_{\RVM}^{\crit}(s)=P_{\RVM}^{\crit}(1-s),
 \qquad
 P_{\RVM}^{\crit}(\bar s)=\overline{P_{\RVM}^{\crit}(s)}.
\end{equation}
Its zeros are exactly
\begin{equation}\label{eq:synthetic-divisor}
 \frac12\pm i\gamma_k^{\RVM},
 \qquad k\ge1,
\end{equation}
and they are simple.  Consequently $\zeta_{\RVM}$ has a simple pole at $s=1$ of residue $1$, the usual trivial zeros at the negative even integers, and nontrivial zeros exactly at \eqref{eq:synthetic-divisor}.  Moreover
$\xi_{\RVM}(s)=\xi_{\RVM}(1-s)$.
\end{theorem}

\begin{proof}
By \eqref{eq:clock-growth},
$\sum_k(\gamma_k^{\RVM})^{-2}<\infty$.  On a compact set $K$,
\[
 \sup_{s\in K}
 \left|\frac{s(s-1)}{1/4+(\gamma_k^{\RVM})^2}\right|
 \ll_K (\gamma_k^{\RVM})^{-2},
\]
so the Weierstrass criterion gives normal convergence.  Each factor depends on $s$ only through $s(s-1)$, which proves the first symmetry; real coefficients give the second.

The $k$th factor vanishes when
\[
 \left(s-\frac12\right)^2+(\gamma_k^{\RVM})^2=0.
\]
The clock is strictly increasing, so the ordinates $\gamma_k^{\RVM}$ are distinct.  The derivative of the vanishing quadratic factor is $2s-1=\pm2i\gamma_k^{\RVM}\ne0$, and no other factor vanishes there.  Hence all nontrivial zeros are simple.

At $s=1$, $P_{\RVM}^{\crit}(1)=1$ and
$G(s)=(s-1)(1+O(s-1))$, so the residue is $1$.  At a negative even integer the pole of $\Gamma(s/2)$ produces a zero of $1/G(s)$; the product $P_{\RVM}^{\crit}$ does not vanish there because $s(s-1)>0$.  The functional equation for $\xi_{\RVM}$ follows from \eqref{eq:P-symmetry}.
\end{proof}

\begin{remark}[Criticality is an input, not RH]\label{rem:not-rh}
The real part $1/2$ in \eqref{eq:synthetic-zero-notation} is fixed by definition.  Theorem~\ref{thm:canonical} is therefore a theorem about a synthetic function.  It gives no information about the genuine real parts $\beta_k$ in \eqref{eq:true-zeros}.
\end{remark}

\section{RVM exponential sums, centroids, and the height transform}

For $u\in\R$ and $N\ge1$ define
\begin{equation}\label{eq:S-RVM}
 S_{\RVM}(N;u):=\sum_{k\le N}e^{iu\gamma_k^{\RVM}},
 \qquad
 C_{\RVM}^{\pts}(N;u):=\frac{S_{\RVM}(N;u)}{N},
\end{equation}
and
\begin{equation}\label{eq:path-centroid}
 C_{\RVM}^{\pathc}(N;u)
 :=\frac1N\sum_{n\le N}S_{\RVM}(n;u)
 =\sum_{k\le N}\frac{N-k+1}{N}e^{iu\gamma_k^{\RVM}}.
\end{equation}

\begin{theorem}[Uniform logarithmic exponential-sum bound]\label{thm:exp-sum}
Let $0<a\le b<\infty$.  Then
\begin{equation}\label{eq:exp-sum-bound}
 \boxed{
 \sup_{u\in[a,b]}|S_{\RVM}(N;u)|\ll_{a,b}1+\log N.}
\end{equation}
\end{theorem}

\begin{proof}
Write $r_k=\gamma_k^{\RVM}$ and $\ell_k=\log(r_k/2\pi)$.  Let $r(x)$ be the inverse of $\mathcal N_{\RVM}$.  From
\[
 r'(x)=\frac{2\pi}{\log(r(x)/2\pi)},
 \qquad
 r''(x)=-\frac{4\pi^2}{r(x)\log^3(r(x)/2\pi)},
\]
we obtain
\begin{equation}\label{eq:rvm-gap-bounds}
 \frac{2\pi}{\ell_{k+1}}
 \le r_{k+1}-r_k
 \le\frac{2\pi}{\ell_k},
\end{equation}
\begin{equation}\label{eq:rvm-second-difference}
 |r_{k+2}-2r_{k+1}+r_k|
 \le\frac{4\pi^2}{r_k\ell_k^3}.
\end{equation}
Choose $k_0=k_0(b)$ so that
$0<\Delta_k:=u(r_{k+1}-r_k)\le\pi$ for $u\in[a,b]$ and $k\ge k_0$.  With $E_k=e^{iur_k}$,
\[
 E_k=\frac{E_{k+1}-E_k}{e^{i\Delta_k}-1}.
\]
Put $b_k=(e^{i\Delta_k}-1)^{-1}$.  Summation by parts gives a boundary term $O_{a,b}(\ell_N)$ and a total-variation term involving $b_{k+1}-b_k$.  Since $|e^{ix}-1|\asymp x$ for $0<x\le\pi$,
\[
 |b_k|\ll_{a,b}\ell_k,
\]
and, using \eqref{eq:rvm-second-difference},
\[
 |b_{k+1}-b_k|
 \ll_{a,b}
 \frac{|\Delta_{k+1}-\Delta_k|}{\Delta_k\Delta_{k+1}}
 \ll_{a,b}\frac1{r_k\ell_k}
 \ll\frac1k.
\]
Here the last comparison follows from the defining clock equation, which gives $r_k\ell_k\asymp k$.  Summing $1/k$ proves \eqref{eq:exp-sum-bound}; the finitely many indices below $k_0$ contribute only a constant.
\end{proof}

\begin{corollary}[Point centroid]\label{cor:point-centroid}
Uniformly for $u\in[a,b]\subset(0,\infty)$,
\begin{equation}
 C_{\RVM}^{\pts}(N;u)=O_{a,b}\!\left(\frac{1+\log N}{N}\right)\longrightarrow0.
\end{equation}
Theorem~\ref{thm:exp-sum} alone does not imply convergence of the path centroid \eqref{eq:path-centroid}.
\end{corollary}

\begin{proof}
The first assertion is immediate from \eqref{eq:exp-sum-bound}.  For the second, the identity
$C_{\RVM}^{\pathc}=N^{-1}\sum_{n\le N}S_{\RVM}(n;u)$ together with the bound $S_{\RVM}(n;u)=O(\log n)$ gives only an $O(\log N)$ bound and therefore no limiting statement.
\end{proof}

Define the midpoint count
\begin{equation}\label{eq:NRVM-count}
 N_{\RVM}(H):=\#\{k:\gamma_k^{\RVM}\le H\}.
\end{equation}
For $x>0$ and $N_{\RVM}(H)\ge1$, define
\begin{equation}\label{eq:MH-RVM}
 M_H^{\RVM}(x)
 :=\frac1{N_{\RVM}(H)}\sum_{\gamma_k^{\RVM}\le H}x^{i\gamma_k^{\RVM}}.
\end{equation}
For the genuine zeros reserve
\begin{equation}\label{eq:MH-genuine}
 M_H(x)
 :=\frac1{N_\zeta(H)}
 \sum_{\substack{\rho=\beta+i\gamma\\0<\gamma\le H}}x^{i\gamma},
\end{equation}
with multiplicities.

\begin{corollary}[Synthetic height-transform decay]\label{cor:MH-decay}
For every fixed $x>1$,
\begin{equation}
 M_H^{\RVM}(x)\longrightarrow0
 \qquad(H\to\infty).
\end{equation}
More precisely, with $N=N_{\RVM}(H)$,
\[
 M_H^{\RVM}(x)=O_x\!\left(\frac{1+\log N}{N}\right).
\]
\end{corollary}

\begin{proof}
Take $u=\log x>0$ in Theorem~\ref{thm:exp-sum}.  Then \eqref{eq:MH-RVM} is exactly the point centroid at $N=N_{\RVM}(H)$.
\end{proof}

\begin{remark}
Neither \eqref{eq:MH-RVM} nor \eqref{eq:MH-genuine} is the Mertens function.  They are finite normalized phase averages.
\end{remark}

\section{The RVM explicit formula, $J_{\RVM}$, and $\pi_{\RVM}$}

For $x>1$, put $u=\log x$ and define
\begin{equation}\label{eq:J-RVM}
 J_{\RVM}(x)
 :=\li(x)
 -2\Re\sum_{k\ge1}\Ei\!\left(\left(\frac12+i\gamma_k^{\RVM}\right)\log x\right)
 -\log2
 +\int_x^\infty\frac{dt}{t(t^2-1)\log t}.
\end{equation}
The zero sum is taken in midpoint-clock order.

\begin{theorem}[Compact-uniform existence of the RVM zero sum]\label{thm:J-exists}
The series
\[
 \sum_{k\ge1}
 \Re\Ei\!\left(\left(\frac12+i\gamma_k^{\RVM}\right)u\right)
\]
converges uniformly for $u$ in every compact interval $[a,b]\subset(0,\infty)$.  Hence $J_{\RVM}(x)$ is an ordinary continuous function on $(1,\infty)$.
\end{theorem}

\begin{proof}
Let $r_k=\gamma_k^{\RVM}$ and
$A_K(u)=\sum_{k\le K}e^{iur_k}$.  By Theorem~\ref{thm:exp-sum},
$A_K(u)=O_{a,b}(1+\log K)$ uniformly on $[a,b]$.  Abel summation therefore gives uniform convergence of
\[
 \sum_{k\ge1}\frac{e^{iur_k}}{r_k}.
\]
Indeed $r_k\asymp k/\log k$ and
\[
 \left|\frac1{r_k}-\frac1{r_{k+1}}\right|
 \ll\frac{r_{k+1}-r_k}{r_k^2}
 \ll\frac1{r_k^2\ell_k},
\]
so the differentiated-weight majorant is $O(\log^2k/k^2)$ after multiplication by $|A_k|$.

Uniformly for $u\in[a,b]$, the standard right-half-plane asymptotic for $\Ei$ gives
\begin{equation}\label{eq:Ei-asymptotic}
 \Re\Ei((1/2+ir_k)u)
 =\frac{e^{u/2}}{u}\frac{\sin(r_ku)}{r_k}
 +O_{a,b}(r_k^{-2}).
\end{equation}
The first series in \eqref{eq:Ei-asymptotic} converges uniformly by the preceding Abel argument and the error series converges absolutely.  This proves the theorem.
\end{proof}

Define the support-truncated version
\begin{equation}\label{eq:J-star}
 J_{\RVM}^{*}(x):=
 \begin{cases}
 J_{\RVM}(x),&x\ge2,\\
 0,&0<x<2.
 \end{cases}
\end{equation}
Set $\pi_{\RVM}(x)=0$ for $0<x<2$, and for $x\ge2$ define
\begin{equation}\label{eq:pi-RVM}
 \boxed{
 \pi_{\RVM}(x)
 :=\sum_{1\le m\le\log_2x}
 \frac{\mu(m)}m J_{\RVM}^{*}(x^{1/m}).}
\end{equation}

\begin{proposition}[Exact finite prime-power inversion]\label{prop:finite-inversion}
For every $x\ge2$,
\begin{equation}\label{eq:J-pi-inverse}
 \boxed{
 J_{\RVM}^{*}(x)
 =\sum_{1\le m\le\log_2x}\frac1m\pi_{\RVM}(x^{1/m}).}
\end{equation}
\end{proposition}

\begin{proof}
All sums are finite because $x^{1/m}<2$ when $m>\log_2x$.  Substitute \eqref{eq:pi-RVM} into the right-hand side of \eqref{eq:J-pi-inverse} and collect the coefficient of $J_{\RVM}^{*}(x^{1/r})$.  It equals
\[
 \frac1r\sum_{d\mid r}\mu(d),
\]
which is $1$ for $r=1$ and $0$ otherwise.  This is ordinary M\"obius inversion in the exponent variable.
\end{proof}

\begin{hypothesis}[Lambert-clock square-root estimate]\label{hyp:Lambert-estimate}
The midpoint explicit formula satisfies
\begin{equation}\label{eq:J-baseline}
 J_{\RVM}(x)-\li(x)
 =O\!\left(\frac{\sqrt x(\log\log x)^2}{\log x}\right).
\end{equation}
Under the same blockwise uniformity, finite M\"obius inversion gives
\begin{equation}\label{eq:pi-baseline}
 \pi_{\RVM}(x)-\li(x)
 =O\!\left(\frac{\sqrt x(\log\log x)^2}{\log x}\right).
\end{equation}
\end{hypothesis}

\begin{remark}[Proof obligation]
Theorem~\ref{thm:J-exists} proves existence and compact-uniform convergence; it does not prove Hypothesis~\ref{hyp:Lambert-estimate}.  A proof of \eqref{eq:J-baseline} must control the branches and symmetric tails of $\Ei$, endpoint terms under integration by parts, and the van der Corput constants uniformly over the required Lambert-clock blocks.  The exponent $1/2$ in the target comes from the inserted synthetic real part $\Re\rho_k^{\RVM}=1/2$ and is not a deduction of classical RH.
\end{remark}

\section{Integer weights and the continuous-to-atomic problem}

Define
\begin{equation}\label{eq:cn}
 c_n^{\RVM}:=\pi_{\RVM}(n)-\pi_{\RVM}(n-1),
 \qquad n\ge2,
\end{equation}
and, initially for $\Re s>1$,
\begin{equation}\label{eq:weighted-zeta}
 \zeta_{\mathbb Z,\RVM}(s)
 :=\prod_{n\ge2}(1-n^{-s})^{-c_n^{\RVM}}.
\end{equation}
This is an analytic weighted Euler object; no positivity or $0/1$ property is built into \eqref{eq:cn}.

For later comparison define the integer-step prime-power count
\begin{equation}\label{eq:JZ}
 J_{\mathbb Z,\RVM}(x)
 :=\sum_{1\le m\le\log_2x}\frac1m
 \pi_{\RVM}(\lfloor x^{1/m}\rfloor).
\end{equation}
Also distinguish the right-continuous extension
\begin{equation}
 J_{\RVM}^{*,\mathbb Z}(x):=J_{\RVM}^{*}(\lfloor x\rfloor)
\end{equation}
from the continuous function $J_{\RVM}^{*}(x)$.

\begin{hypothesis}[Uniform unit-cell control]\label{hyp:unit-cell}
For sufficiently large $x$ and $0<h\le1$,
\begin{equation}\label{eq:unit-cell}
 |J_{\RVM}(x+h)-J_{\RVM}(x)|
 \ll h+\frac{\log x}{\sqrt x},
\end{equation}
uniformly in the unit interval.  In particular, finite inversion then yields
\begin{equation}
 \pi_{\RVM}(y)-\pi_{\RVM}(\lfloor y\rfloor)=O(1),
 \qquad c_n^{\RVM}=O(1).
\end{equation}
\end{hypothesis}

\begin{hypothesis}[Mellin--Stieltjes identification]\label{hyp:MS}
The logarithm of the weighted Euler product \eqref{eq:weighted-zeta} and the continuous explicit-formula object are related by the standard Mellin--Stieltjes transform of their prime-power counts on a common half-plane, with the continuous unit-cell defect represented by an ordinary holomorphic integral.
\end{hypothesis}

\begin{corollary}[Conditional weighted factorization]\label{cor:weighted-factorization}
Under Hypotheses~\ref{hyp:unit-cell} and \ref{hyp:MS},
\begin{equation}\label{eq:JZ-diff}
 J_{\mathbb Z,\RVM}(x)-J_{\RVM}^{*}(x)=O(\log\log x),
\end{equation}
and
\begin{equation}\label{eq:weighted-factorization}
 \boxed{
 \zeta_{\mathbb Z,\RVM}(s)=\zeta_{\RVM}(s)e^{h_{\RVM}(s)},}
\end{equation}
where
\begin{equation}\label{eq:hRVM}
 h_{\RVM}(s)
 :=s\int_1^\infty
 \bigl(J_{\mathbb Z,\RVM}(x)-J_{\RVM}^{*}(x)\bigr)x^{-s-1}\,dx
\end{equation}
is holomorphic for $\Re s>0$.  Thus the two objects have the same zero divisor in $\Re s>0$.
\end{corollary}

\begin{proof}
Under \eqref{eq:unit-cell}, replacing each $x^{1/m}$ by its integer part costs $O(1/m)$ in the $m$th summand.  Summing $m^{-1}$ up to $\log_2x$ gives \eqref{eq:JZ-diff}.  Hypothesis~\ref{hyp:MS} then turns the prime-power-count difference into \eqref{eq:hRVM}; logarithmic growth makes the integral locally normally convergent for $\Re s>0$.  Exponentiation introduces no zeros.
\end{proof}

The first Poisson harmonic predicted by the midpoint clock is the unit-periodic term $2\cos(2\pi x)$.  Its integer-cell contribution cancels exactly:
\begin{lemma}[Exact first-harmonic cell cancellation]\label{lem:cell-cancel}
For every integer $n$,
\begin{equation}\label{eq:cell-cancel}
 \boxed{\int_{n-1}^{n}2\cos(2\pi x)\,dx=0.}
\end{equation}
\end{lemma}

\begin{hypothesis}[Uniform cell remainder]\label{hyp:cell-remainder}
The increment after removing the main density and the first harmonic satisfies
\begin{equation}\label{eq:cell-remainder}
 c_n^{\RVM}
 -\int_{n-1}^{n}\left(\frac1{\log x}+2\cos(2\pi x)\right)dx
 =o\!\left(\frac1{\log n}\right).
\end{equation}
\end{hypothesis}

\begin{proposition}[Conditional eventual positivity]\label{prop:eventual-positive}
Hypothesis~\ref{hyp:cell-remainder} implies
\begin{equation}
 c_n^{\RVM}\sim\frac1{\log n},
\end{equation}
and hence $0<c_n^{\RVM}<1$ for all sufficiently large $n$.
\end{proposition}

\begin{proof}
The main-density integral is
$1/\log n+O(1/(n\log^2n))$.  The first harmonic vanishes by Lemma~\ref{lem:cell-cancel}, and the assumed remainder is $o(1/\log n)$.
\end{proof}

The passage from these analytic weights to an ordinary $0/1$ set of integer ``RVM primes'' is therefore conditional.  It must not be confused with a theorem about the genuine primes.

\section{The critical RVM Mertens transform}

Define
\begin{equation}\label{eq:Mcrit-transform}
 \widehat M_{\RVM}^{\crit}(s):=\frac1{\zeta_{\RVM}(s)}.
\end{equation}
Whenever ordinary inversion is justified, put
\begin{equation}\label{eq:Mcrit-inverse}
 M_{\RVM}^{\crit}(x)
 :=\frac1{2\pi i}\int_{c-i\infty}^{c+i\infty}
 \frac{x^s}{s\zeta_{\RVM}(s)}\,ds,
 \qquad c>1.
\end{equation}
If the integral is not justified as an ordinary inverse, \eqref{eq:Mcrit-inverse} is retained only as a formal inverse-Mellin notation until a function or distribution space is specified.

\begin{proposition}[Synthetic pole abscissa]\label{prop:Mcrit-poles}
The nontrivial poles of $\widehat M_{\RVM}^{\crit}$ are exactly
\begin{equation}
 \frac12\pm i\gamma_k^{\RVM}.
\end{equation}
Thus the nontrivial pole abscissa of the transform is exactly $1/2$ by construction.
\end{proposition}

\begin{proof}
This is immediate from Theorem~\ref{thm:canonical} and \eqref{eq:Mcrit-transform}.
\end{proof}

\begin{proposition}[Gaussian square-root bound under vertical growth]\label{prop:Gaussian}
Fix $c>1$ and $0<\eps<c-1/2$.  Assume uniformly for
$1/2+\eps\le\sigma\le c$ that
\begin{equation}\label{eq:subgaussian}
 \frac1{\zeta_{\RVM}(\sigma+it)}=\exp(o(t^2))
 \qquad(|t|\to\infty),
\end{equation}
with the same control on the horizontal sides of the contour rectangles.  Define
\begin{equation}\label{eq:MG}
 M_{\RVM,G}^{\crit}(x)
 :=\frac1{2\pi i}\int_{c-i\infty}^{c+i\infty}
 e^{s^2}\frac{x^s}{s\zeta_{\RVM}(s)}\,ds.
\end{equation}
Then
\begin{equation}\label{eq:MG-bound}
 M_{\RVM,G}^{\crit}(x)=O_\eps(x^{1/2+\eps}).
\end{equation}
\end{proposition}

\begin{proof}
On $s=\sigma+it$, $|e^{s^2}|=e^{\sigma^2-t^2}$.  Assumption \eqref{eq:subgaussian} gives absolute convergence on the vertical lines and makes the horizontal integrals tend to zero.  The strip $1/2+\eps\le\Re s\le c$ contains no nontrivial pole of $1/\zeta_{\RVM}$; $s=1$ is a zero, not a pole, of that reciprocal.  Cauchy's theorem therefore shifts the contour in \eqref{eq:MG} to $\Re s=1/2+\eps$.  The factor $x^s$ then contributes $x^{1/2+\eps}$ and the remaining integral is finite by the same Gaussian domination.
\end{proof}

\begin{remark}[Why the transition theorem is needed]\label{rem:unsmoothed-separate}
The location of the synthetic poles alone does not imply an unsmoothed square-root bound.  The required vertical estimate is supplied later by Theorem~\ref{thm:transition}, which also fixes the symmetric interpretation of the Perron integral and proves Theorem~\ref{thm:unsmoothed-RVM}.  Thus no unsmoothed estimate is inferred merely from the fact that the synthetic divisor lies on $\Re s=1/2$.
\end{remark}

\section{Exact comparison with the classical zeta, primes, and Mertens transform}

Put
\begin{equation}\label{eq:Pzeta}
 P_\zeta(s):=G(s)\zeta(s)=2\xi(s).
\end{equation}

\begin{theorem}[Exact paired Hadamard normalization]\label{thm:Hadamard}
The paired product
\begin{equation}\label{eq:paired-product}
 P_\zeta^{\rm pair}(s)
 :=\prod_{k\ge1}
 \left(1-\frac{s}{\rho_k}\right)
 \left(1-\frac{s}{\overline{\rho_k}}\right)
\end{equation}
converges absolutely and locally uniformly and satisfies
\begin{equation}\label{eq:Hadamard-exact}
 \boxed{P_\zeta(s)=P_\zeta^{\rm pair}(s).}
\end{equation}
In particular, there is no residual factor $e^{a+bs}$.
\end{theorem}

\begin{proof}
For $\rho=\beta+i\gamma$ with $0<\beta<1$,
\[
 \left(1-\frac{s}{\rho}\right)
 \left(1-\frac{s}{\bar\rho}\right)
 =1+O_K(\gamma^{-2})
\]
uniformly on compact sets.  Since $N_\zeta(T)=O(T\log T)$,
$\sum_{\gamma>0}\gamma^{-2}<\infty$, proving normal convergence.  Functional symmetry pairs each off-line zero $\beta+i\gamma$ with $1-\beta+i\gamma$ at the same ordinate; grouping these absolutely convergent factors gives
$P_\zeta^{\rm pair}(1-s)=P_\zeta^{\rm pair}(s)$.  The product has order at most one, has exactly the zeros of $P_\zeta$ with multiplicity, and hence Hadamard factorization gives
$P_\zeta=e^{a+bs}P_\zeta^{\rm pair}$.  Both functions equal $1$ at $s=0$, so $e^a=1$.  Applying $s\mapsto1-s$ and using the two functional symmetries yields $e^{bs}=e^{b(1-s)}$ for every $s$, hence $b=0$.  The remaining constant is $1$.
\end{proof}

Define the exact quotient
\begin{equation}\label{eq:Q-RVM}
 \boxed{
 \mathcal Q_{\RVM}(s)
 :=\frac{\zeta(s)}{\zeta_{\RVM}(s)}
 =\frac{P_\zeta(s)}{P_{\RVM}^{\crit}(s)}.}
\end{equation}

\begin{proposition}[Exact transform relation]\label{prop:M-transform-relation}
For $\Re s>1$, and by meromorphic continuation wherever both sides are defined,
\begin{equation}\label{eq:M-transform-relation}
 \boxed{
 \widehat M_{\RVM}^{\crit}(s)
 =\mathcal Q_{\RVM}(s)\frac1{\zeta(s)}.}
\end{equation}
If the inverse Mellin transform of $\mathcal Q_{\RVM}$ and the inverse defining $M_{\RVM}^{\crit}$ exist in a common distribution space, with Mellin convolution defined there, then
\begin{equation}\label{eq:M-convolution}
 M_{\RVM}^{\crit}
 =\kappa_{\RVM}^{\crit}\Mell M.
\end{equation}
No boundedness or invertibility of $\mathcal Q_{\RVM}$ is asserted.
\end{proposition}

\begin{proof}
Equation \eqref{eq:M-transform-relation} is the algebraic identity
$1/\zeta_{\RVM}=(\zeta/\zeta_{\RVM})(1/\zeta)$.  The convolution statement is exactly the Mellin convolution theorem under the stated existence assumptions.
\end{proof}

On the prime side define
\begin{equation}\label{eq:prime-defect}
 \eps_n^{\RVM}
 :=\mathbf 1_{\{n\ \mathrm{prime}\}}-c_n^{\RVM}.
\end{equation}
Then for every integer $N\ge2$,
\begin{equation}\label{eq:prime-defect-cumulative}
 \boxed{
 \sum_{2\le n\le N}\eps_n^{\RVM}
 =\pi(N)-\pi_{\RVM}(N).}
\end{equation}
This is an exact identity, not an asymptotic statement.

\section{Continuous transport from the RVM clock to the genuine ordinates}

This section incorporates the three-step transport construction.  To keep the orientation used in the present paper, we parameterize the path from the synthetic clock to the genuine ordinates by $\kappa\in[0,1]$:
\begin{equation}\label{eq:transport-delta}
 \delta_k:=\gamma_k-\gamma_k^{\RVM},
 \qquad
 \gamma_{k,\kappa}
 :=\gamma_k^{\RVM}+\kappa\delta_k.
\end{equation}
Thus $\kappa=0$ is the RVM endpoint and $\kappa=1$ is the list of genuine ordinates projected onto the critical line.  The parameter $\lambda$ used in the three-step note is related by $\lambda=1-\kappa$.

Define
\begin{equation}\label{eq:path-product}
 P_\kappa(s)
 :=\prod_{k\ge1}
 \left(1+\frac{s(s-1)}{1/4+\gamma_{k,\kappa}^2}\right),
 \qquad
 \zeta_\kappa(s):=\frac{P_\kappa(s)}{G(s)}.
\end{equation}
Also put
\begin{equation}\label{eq:Pline}
 P_{\linez}(s):=P_1(s),
 \qquad
 \zeta_{\linez}(s):=\zeta_1(s).
\end{equation}
The object denoted $P_{Z1}$ or $\zeta_{Z1}$ in the three-step note is precisely the endpoint $P_{\linez}$ or $\zeta_{\linez}$ here.

\begin{theorem}[Joint normal convergence of the transport family]\label{thm:path}
The product \eqref{eq:path-product} converges locally uniformly jointly on $[0,1]\times\C$.  For every $\kappa$, $P_\kappa$ is entire and
\begin{equation}
 P_\kappa(1-s)=P_\kappa(s),
 \qquad
 P_\kappa(\bar s)=\overline{P_\kappa(s)}.
\end{equation}
Its zeros, with multiplicity, are exactly
\begin{equation}\label{eq:path-divisor}
 \frac12\pm i\gamma_{k,\kappa}.
\end{equation}
For $0\le\kappa<1$ the positive ordinates are strictly increasing and the zeros in \eqref{eq:path-divisor} are simple.  At $\kappa=1$ their multiplicities are exactly the multiplicities of the genuine ordinates.  Moreover
\begin{equation}\label{eq:path-endpoints}
 \boxed{
 \zeta_0=\zeta_{\RVM},
 \qquad
 \zeta_1=\zeta_{\linez}.}
\end{equation}
For every integer $j\ge0$, $P_\kappa^{(j)}\to P_{\kappa_0}^{(j)}$ locally uniformly as $\kappa\to\kappa_0$.
\end{theorem}

\begin{proof}
Lemma~\ref{lem:bounded-displacement} gives $|\delta_k|\le D$ after enlarging $D$ to absorb finitely many initial indices.  Uniformly in $\kappa$,
$\gamma_{k,\kappa}\asymp k/\log k$, and for all large $k$ it is bounded below by a positive constant multiple of $\gamma_k^{\RVM}$.  Therefore, on every compact $K\subset\C$,
\[
 \sum_{k\ge1}
 \sup_{\substack{s\in K\\0\le\kappa\le1}}
 \left|\frac{s(s-1)}{1/4+\gamma_{k,\kappa}^2}\right|
 <\infty.
\]
This proves joint normal convergence.  The symmetry and the divisor follow factor by factor.

For $\kappa<1$,
\[
 \gamma_{k+1,\kappa}-\gamma_{k,\kappa}
 =(1-\kappa)(\gamma_{k+1}^{\RVM}-\gamma_k^{\RVM})
 +\kappa(\gamma_{k+1}-\gamma_k)>0,
\]
because the RVM clock is strictly increasing and the genuine ordinate list is nondecreasing.  At $\kappa=1$, equality occurs exactly when genuine ordinates are repeated.  The endpoint identities are definitions.  Local uniform convergence of all $s$-derivatives follows from Cauchy's integral formula.
\end{proof}

\begin{theorem}[Infinitesimal transport generator]\label{thm:generator}
Let $q(s)=s(s-1)$.  Off the moving divisor of $P_\kappa$,
\begin{equation}\label{eq:generator}
 \boxed{
 \partial_\kappa\log P_\kappa(s)
 =-2q(s)\sum_{k\ge1}
 \frac{\gamma_{k,\kappa}\delta_k}
 {(1/4+\gamma_{k,\kappa}^2)
  (1/4+\gamma_{k,\kappa}^2+q(s))}.}
\end{equation}
The series converges locally uniformly on compact subsets separated from the moving divisor.  Since $G$ is independent of $\kappa$,
\begin{equation}\label{eq:generator-zeta}
 \partial_\kappa\log\zeta_\kappa(s)
 =\partial_\kappa\log P_\kappa(s),
 \qquad
 \partial_\kappa\log\frac1{\zeta_\kappa(s)}
 =-\partial_\kappa\log P_\kappa(s).
\end{equation}
\end{theorem}

\begin{proof}
Set $A_{k,\kappa}=1/4+\gamma_{k,\kappa}^2$.  Since
$\partial_\kappa A_{k,\kappa}=2\gamma_{k,\kappa}\delta_k$,
\[
 \partial_\kappa
 \log\left(1+\frac{q}{A_{k,\kappa}}\right)
 =-\frac{2q\gamma_{k,\kappa}\delta_k}
 {A_{k,\kappa}(A_{k,\kappa}+q)}.
\]
On compact sets separated from the divisor the summand is
$O_K(|\delta_k|/\gamma_k^3)$.  Bounded displacement and
$\gamma_k\asymp k/\log k$ imply
$\sum_k|\delta_k|/\gamma_k^3<\infty$, so termwise differentiation is justified.
\end{proof}

Write the right-hand side of \eqref{eq:generator} as $\mathcal V_\kappa(s)$.  Integration gives an exact transport formula.

\begin{proposition}[Exact integrated multiplier]\label{prop:integrated-multiplier}
On every component avoiding the moving divisor,
\begin{equation}\label{eq:integrated-zeta}
 \boxed{
 \zeta_\kappa(s)
 =\zeta_{\RVM}(s)
 \exp\left(\int_0^\kappa\mathcal V_u(s)\,du\right),}
\end{equation}
\begin{equation}\label{eq:integrated-inverse}
 \boxed{
 \frac1{\zeta_\kappa(s)}
 =\frac1{\zeta_{\RVM}(s)}
 \exp\left(-\int_0^\kappa\mathcal V_u(s)\,du\right).}
\end{equation}
At the endpoint,
\begin{equation}\label{eq:integrated-product-ratio}
 \exp\left(\int_0^1\mathcal V_u(s)\,du\right)
 =\prod_{k\ge1}
 \frac{1+q(s)/(1/4+\gamma_k^2)}
      {1+q(s)/(1/4+(\gamma_k^{\RVM})^2)}
 =\frac{P_{\linez}(s)}{P_{\RVM}^{\crit}(s)}.
\end{equation}
The exponential identities extend meromorphically across the plane by the product identity.
\end{proposition}

\begin{proof}
The first two identities are the fundamental theorem of calculus applied to Theorem~\ref{thm:generator}.  For the last identity, integrate each mode using $r=\gamma_{k,u}$ and $dr=\delta_k\,du$:
\[
 \int_0^1
 \frac{-2q\,\gamma_{k,u}\delta_k}
 {(1/4+\gamma_{k,u}^2)(1/4+\gamma_{k,u}^2+q)}\,du
 =\log\frac{1+q/(1/4+\gamma_k^2)}
 {1+q/(1/4+(\gamma_k^{\RVM})^2)}.
\]
Normal convergence permits summation.  Exponentiating gives \eqref{eq:integrated-product-ratio} and removes any branch ambiguity.
\end{proof}

\begin{corollary}[First-order endpoint expansion]\label{cor:endpoint-expansion}
Let $K\subset\C$ be compact and disjoint from the zero divisor of $P_{\linez}$.  Uniformly for $s\in K$ as $\kappa\uparrow1$,
\begin{equation}\label{eq:endpoint-zeta-expansion}
 \zeta_\kappa(s)=\zeta_{\linez}(s)
 \left(1-(1-\kappa)\mathcal V_1(s)+O_K((1-\kappa)^2)\right),
\end{equation}
and
\begin{equation}\label{eq:endpoint-inverse-expansion}
 \frac1{\zeta_\kappa(s)}=\frac1{\zeta_{\linez}(s)}
 \left(1+(1-\kappa)\mathcal V_1(s)+O_K((1-\kappa)^2)\right).
\end{equation}
\end{corollary}

\begin{proof}
On $K$ the series defining $\mathcal V_\kappa$ and its $\kappa$-derivative converge normally for $\kappa$ sufficiently close to $1$, because $|\delta_k|=O(1)$ and the differentiated summands are $O_K(\gamma_k^{-4})+O_K(\gamma_k^{-3})$.  Hence
$\mathcal V_\kappa=\mathcal V_1+O_K(1-\kappa)$.  From
$\log(\zeta_\kappa/\zeta_{\linez})=-\int_\kappa^1\mathcal V_u\,du$ we obtain
$\log(\zeta_\kappa/\zeta_{\linez})=-(1-\kappa)\mathcal V_1+O_K((1-\kappa)^2)$; exponentiation gives \eqref{eq:endpoint-zeta-expansion}, and inversion gives \eqref{eq:endpoint-inverse-expansion}.
\end{proof}

\begin{proposition}[Transport of simple residues]\label{prop:residues}
Suppose the genuine ordinate $\gamma_k$ is not repeated.  Put
$\rho_{k,\kappa}=1/2+i\gamma_{k,\kappa}$.  For $\kappa<1$ define
\begin{equation}\label{eq:residue-transport}
 B_{k,\kappa}
 :=\Res_{s=\rho_{k,\kappa}}
 \frac{1}{s\zeta_\kappa(s)}
 =\frac{G(\rho_{k,\kappa})}
 {\rho_{k,\kappa}P_\kappa'(\rho_{k,\kappa})}.
\end{equation}
Then
\begin{equation}
 B_{k,\kappa}\longrightarrow
 \frac{G(1/2+i\gamma_k)}
 {(1/2+i\gamma_k)P_{\linez}'(1/2+i\gamma_k)}
 \qquad(\kappa\uparrow1).
\end{equation}
Consequently each corresponding simple residue mode converges pointwise in $x>0$.
\end{proposition}

\begin{proof}
Theorem~\ref{thm:path} gives joint convergence of the zero locations and local uniform convergence of the first derivatives.  The denominator in \eqref{eq:residue-transport} has a nonzero limit because the endpoint zero is simple by hypothesis.  Continuity of $G$ finishes the proof.
\end{proof}

If the genuine ordinate has multiplicity greater than one, simple residues coalesce and one must transport the full principal part rather than apply Proposition~\ref{prop:residues} term by term.

\section{The endpoint criterion and the meaning of the three-step construction}

The line endpoint \eqref{eq:Pline} depends on the genuine ordinates $\gamma_k$ but discards the genuine real parts $\beta_k$.  It therefore exists unconditionally, regardless of RH.

\begin{theorem}[Endpoint identification]\label{thm:endpoint-RH}
\begin{equation}\label{eq:endpoint-RH}
 \boxed{
 \zeta_{\linez}(s)\equiv\zeta(s)
 \quad\Longleftrightarrow\quad
 \mathrm{RH}.}
\end{equation}
\end{theorem}

\begin{proof}
If $\zeta_{\linez}=\zeta$, the nontrivial zero divisor of $\zeta$ is the divisor of $\zeta_{\linez}$, hence all nontrivial zeros lie on $\Re s=1/2$.

Conversely assume RH.  Then every genuine upper-half-plane zero is $1/2+i\gamma_k$, with multiplicity.  The exact paired Hadamard product of Theorem~\ref{thm:Hadamard} therefore becomes
\[
 P_\zeta(s)
 =\prod_{k\ge1}
 \left(1+\frac{s(s-1)}{1/4+\gamma_k^2}\right)
 =P_{\linez}(s).
\]
Dividing both sides by the common factor $G(s)$ gives
$\zeta=\zeta_{\linez}$.
\end{proof}

\begin{remark}
Constructing the path $\zeta_{\RVM}\to\zeta_{\linez}$ does not prove RH.  The endpoint exists whether or not RH is true.  What is equivalent to RH is the additional identification of that endpoint with the classical zeta function.
\end{remark}

For the inverse functions put
\begin{equation}\label{eq:Fkappa}
 F_\kappa(s):=\frac1{\zeta_\kappa(s)}.
\end{equation}
For $c>1$ and finite $U$, define the truncated Perron integral
\begin{equation}\label{eq:truncated-Perron}
 M_{\kappa,U}(x)
 :=\frac1{2\pi i}\int_{c-iU}^{c+iU}
 F_\kappa(s)\frac{x^s}{s}\,ds.
\end{equation}

\begin{proposition}[Finite-height Perron transport]\label{prop:finite-Perron}
For every fixed $c>1$, $U<\infty$, and $x>0$,
\begin{equation}
 M_{\kappa,U}(x)\longrightarrow M_{\linez,U}(x)
 \qquad(\kappa\uparrow1),
\end{equation}
where $M_{\linez,U}$ is obtained from \eqref{eq:truncated-Perron} by replacing $F_\kappa$ with $F_{\linez}=1/\zeta_{\linez}$.
\end{proposition}

\begin{proof}
The compact vertical segment $\{c+it:|t|\le U\}$ is disjoint from the nontrivial divisor, and $P_\kappa\to P_{\linez}$ uniformly there.  Hence $F_\kappa\to F_{\linez}$ uniformly on the segment, so the finite integral converges by uniform convergence.
\end{proof}

The passage $U\to\infty$ is a separate analytic operation.  The following theorem is therefore stated with its required hypotheses visible.

\begin{theorem}[Conditional coefficient-level endpoint criterion]\label{thm:coefficient-criterion}
Assume that in $\Re s>1$ the endpoint reciprocal has an absolutely convergent Dirichlet series
\begin{equation}\label{eq:line-Dirichlet}
 F_{\linez}(s)=\frac1{\zeta_{\linez}(s)}
 =\sum_{n\ge1}\frac{\mu_{\linez}(n)}{n^s},
\end{equation}
and that Perron inversion for this series defines a summatory function $M_{\linez}$ with
\begin{equation}\label{eq:line-increments}
 \mu_{\linez}(n)=M_{\linez}(n)-M_{\linez}(n-1)
\end{equation}
under a fixed endpoint convention.  Then the following are equivalent:
\begin{enumerate}[label=\textup{(\roman*)}]
\item $M_{\linez}(n)=M(n)$ for every integer $n\ge1$;
\item $\mu_{\linez}(n)=\mu(n)$ for every $n\ge1$;
\item $1/\zeta_{\linez}(s)=1/\zeta(s)$ for $\Re s>1$;
\item $\zeta_{\linez}(s)\equiv\zeta(s)$;
\item RH is true.
\end{enumerate}
\end{theorem}

\begin{proof}
The equivalence of (i) and (ii) follows from taking discrete increments, with the normalization at $n=0$.  Under absolute convergence, (ii) gives (iii) coefficient by coefficient.  In $\Re s>1$ both zeta functions are nonzero, so (iii) is equivalent to equality there; meromorphic continuation then gives (iv).  The equivalence of (iv) and (v) is Theorem~\ref{thm:endpoint-RH}.  Finally, under (iv), \eqref{eq:line-Dirichlet} is the classical expansion of $1/\zeta$, and Perron inversion gives (i).
\end{proof}

The transport multiplier in inverse form is
\begin{equation}\label{eq:Ktransport}
 \mathcal K_{\rm tr}(s)
 :=\exp\left(-\int_0^1\mathcal V_\kappa(s)\,d\kappa\right)
 =\frac{P_{\RVM}^{\crit}(s)}{P_{\linez}(s)},
\end{equation}
so that
\begin{equation}\label{eq:Fline-factor}
 \boxed{
 F_{\linez}(s)=F_{\RVM}(s)\mathcal K_{\rm tr}(s),
 \qquad F_{\RVM}=\frac1{\zeta_{\RVM}}.}
\end{equation}

\begin{remark}[The coefficient closure proposed in the three-step note]\label{rem:word-coefficient-route}
The Word note proposed expanding both factors on the right of \eqref{eq:Fline-factor} as absolutely convergent ordinary Dirichlet series and then identifying their Dirichlet convolution with $\mu$.  Theorem~\ref{thm:no-Dirichlet-RVM} below proves that the first of those expansions does not exist in any right half-plane.  Thus the proposed convolution is not an available coefficient test in the ordinary Dirichlet-series category.  The endpoint criterion of Theorem~\ref{thm:coefficient-criterion} remains correct under its own explicit endpoint hypotheses, but it cannot be obtained by assigning ordinary coefficients $a_{\RVM}(n)$ to $1/\zeta_{\RVM}$.
\end{remark}

\section{Vertical transition theorem and the unsmoothed synthetic Mertens bound}
\label{sec:transition}

The midpoint product admits a direct vertical evaluation.  This section records the full argument because it changes the status of two points left open in the initial extraction from v56: the unsmoothed inverse of $1/\zeta_{\RVM}$ and the coefficient factorization suggested by the three-step transport note.

Let
\begin{equation}\label{eq:chi}
 \chi(s)
 :=2^s\pi^{s-1}\sin\!\left(\frac{\pi s}{2}\right)\Gamma(1-s)
 =\pi^{s-1/2}\frac{\Gamma((1-s)/2)}{\Gamma(s/2)}.
\end{equation}
For $s=\tfrac12+z$, $z=\sigma'+it$, put
\begin{equation}\label{eq:transition-bw}
 b:=-iz=t-i\sigma',\qquad
 w(s):=\exp\{-2\pi i\,\mathcal N_{\RVM}(b)\}.
\end{equation}
Stirling's formula gives, uniformly on compact $\sigma$-sets,
\begin{equation}\label{eq:chi-size}
 |\chi(\sigma+it)|
 =\left(\frac{|t|}{2\pi}\right)^{1/2-\sigma}(1+O(|t|^{-1})).
\end{equation}

\begin{theorem}[Transition evaluation of the midpoint product]\label{thm:transition}
There is a constant $Z_{\RVM}>0$ such that the following hold.
\begin{enumerate}[label=\textup{(\alph*)}]
\item For every fixed $B>0$ there are $t_0(B)$ and $C_B$ such that, uniformly for $0<\sigma'\le B$ and $t\ge t_0(B)$,
\begin{equation}\label{eq:transition-w}
 \boxed{\zeta_{\RVM}(s)=Z_{\RVM}(1+w(s))(1+E_0(s)),\qquad |E_0(s)|\le C_B/t.}
\end{equation}
\item For every $0<a_0<b_0<1/2$ there are $t_0$ and $C$ such that, uniformly for $a_0\le\sigma'\le b_0$ and $|t|\ge t_0$,
\begin{equation}\label{eq:transition-chi}
 \boxed{\zeta_{\RVM}(s)=Z_{\RVM}(1+\chi(s))(1+E(s)),\qquad |E(s)|\le C/|t|.}
\end{equation}
\end{enumerate}
The statements for $t<0$ follow by conjugation.
\end{theorem}

\begin{proof}
It is enough to take $t>0$.  Write $\gamma_k^{\RVM}=\tau_k$ temporarily inside this proof only, to reduce typographical load; this local abbreviation is not used elsewhere in the paper.  Since $s(s-1)=z^2-1/4$,
\begin{equation}\label{eq:transition-product}
 \frac{P_{\RVM}^{\crit}(s)}{P_{\RVM}^{\crit}(1/2)}
 =\prod_{k\ge1}\left(1+\frac{z^2}{\tau_k^2}\right)
 =\exp\Lambda(z),\qquad
 \Lambda(z)=\sum_{k\ge1}f(\tau_k),
\end{equation}
where $f(\tau)=\Log(1+z^2/\tau^2)$.  For $\tau>0$,
$\Im(1+z^2/\tau^2)=2\sigma't/\tau^2>0$, so the principal logarithm is smooth and
\begin{equation}\label{eq:fprime}
 f'(\tau)=-\frac{2z^2}{\tau(\tau^2+z^2)}
 =-\frac2\tau+\frac1{\tau-iz}+\frac1{\tau+iz}.
\end{equation}

Set $\tau_0=4\pi$.  The counting function of the midpoint ordinates on $[\tau_0,\infty)$ is
\[
 n(\tau)=\left\lfloor\mathcal N_{\RVM}(\tau)+\tfrac12\right\rfloor
 =\mathcal N_{\RVM}(\tau)-\psi(\mathcal N_{\RVM}(\tau)+\tfrac12),
 \qquad \psi(y)=y-\lfloor y\rfloor-\tfrac12.
\]
Stieltjes integration and integration by parts in the sawtooth term give
\begin{equation}\label{eq:transition-split}
 \Lambda(z)=J(z)+\Psi(z),
\end{equation}
with
\begin{align}
 J(z)&=\int_{\tau_0}^{\infty}f\,d\mathcal N_{\RVM}+f(\tau_0)\mathcal N_{\RVM}(\tau_0),\label{eq:Jtransition}\\
 \Psi(z)&=\int_{\tau_0}^{\infty}\psi(\mathcal N_{\RVM}(\tau)+\tfrac12)f'(\tau)\,d\tau.\label{eq:Psitransition}
\end{align}

For the smooth part, a second integration by parts yields
\[
 J(z)=\int_{\tau_0}^{\infty}\mathcal N_{\RVM}(\tau)
 \left(\frac2\tau-\frac{2\tau}{\tau^2+z^2}\right)d\tau.
\]
The constant $7/8$ contributes
$\frac78\Log(1+z^2/\tau_0^2)=\frac74\Log z-\frac74\log\tau_0+O_B(t^{-2})$.
For the remaining term put $c=2\pi e$ and
\[
 I(z):=\frac1\pi\int_0^\infty\log\frac\tau c\,\frac{z^2}{\tau^2+z^2}\,d\tau.
\]
It is holomorphic on $\Re z>0$.  For real $z=x>0$, the substitution $\tau=xu$ and
$\int_0^\infty(1+u^2)^{-1}du=\pi/2$, $\int_0^\infty\log u\,(1+u^2)^{-1}du=0$
give $I(x)=\frac x2\log(x/c)$; hence by the identity theorem
$I(z)=\frac z2\Log(z/c)$.  Removing $[0,\tau_0]$ changes this by a constant plus $O_B(t^{-2})$.  Therefore
\begin{equation}\label{eq:Jasymp}
 J(z)=\frac z2\Log\frac z{2\pi e}+\frac74\Log z+c_1+O_B(t^{-2}).
\end{equation}
On the other hand Stirling's formula, uniformly for $0<\sigma'\le B$, gives
\begin{equation}\label{eq:Gasymp}
 \log G(s)=\frac z2\Log\frac z{2\pi e}+\frac74\Log z+c_2+O_B(t^{-1}),
\end{equation}
so
\begin{equation}\label{eq:smooth-difference}
 J(z)-\log G(s)=c_3+O_B(t^{-1}),\qquad c_3=c_1-c_2\in\R.
\end{equation}

We next treat the sawtooth term.  The clock is holomorphic on $\Re\tau>0$ and, for $\Re\tau\ge\tau_0$,
\begin{equation}\label{eq:clock-halfplane-bound}
 2\pi\Im\mathcal N_{\RVM}(\tau)\ge(\log2)\Im\tau\quad(\Im\tau\ge0),
 \qquad
 2\pi\Im\mathcal N_{\RVM}(\tau)\le(\log2)\Im\tau\quad(\Im\tau\le0).
\end{equation}
Indeed, for $\tau=x+iy$, $x\ge4\pi$, $u=y/x>0$,
\[
 2\pi\Im\mathcal N_{\RVM}(\tau)
 =y\left[\log\frac{x}{2\pi e}+\frac12\log(1+u^2)+\frac{\arctan u}{u}\right]
 \ge y\log2,
\]
because $h(u)=\frac12\log(1+u^2)+\arctan(u)/u-1$ satisfies $h(0^+)=0$ and
$h'(u)=(u-\arctan u)/u^2\ge0$; the lower half-plane follows by conjugation.
In particular $|w(s)|<1$ for $\sigma'>0$, and the vertical rays carry exponential decay $2^{-my}$ in the $m$th Fourier mode.

Regularize the sawtooth by
\[
 \psi_r(y)=-\sum_{m\ge1}r^m\frac{\sin2\pi my}{\pi m},\qquad 0<r<1.
\]
Dominated convergence gives $\Psi=\lim_{r\uparrow1}\Psi_r$, and Fubini is valid for fixed $r$:
\[
 \Psi_r(z)=-\sum_{m\ge1}\frac{(-r)^m}{2\pi i\,m}(I_m^+-I_m^-),
 \qquad
 I_m^\pm=\int_{\tau_0}^\infty e^{\pm2\pi im\mathcal N_{\RVM}(\tau)}f'(\tau)\,d\tau.
\]
Close $I_m^+$ in the upper right quadrant and $I_m^-$ in the lower right quadrant.  The quarter-circle contributions tend to zero by $f'(\tau)=O(|z|^2|\tau|^{-3})$.  The only enclosed moving pole is
$b=-iz=t-i\sigma'$ in the lower quadrant, with residue $1$.  With
\[
 R_m^+=i\int_0^\infty e^{2\pi im\mathcal N_{\RVM}(\tau_0+iy)}f'(\tau_0+iy)\,dy,
\]
\[
 R_m^-=-i\int_0^\infty e^{-2\pi im\mathcal N_{\RVM}(\tau_0-iy)}f'(\tau_0-iy)\,dy,
\]
the orientation gives $I_m^+=R_m^+$ and
$I_m^-=R_m^--2\pi i\,w(s)^m$.  Consequently
\begin{equation}\label{eq:Psi-r}
 \Psi_r(z)=\Log(1+r w(s))+V_r(z),
 \qquad
 V_r(z)=-\sum_{m\ge1}\frac{(-r)^m}{2\pi i\,m}(R_m^+-R_m^-).
\end{equation}
Split $f'=h_0+g_z$ with $h_0(\tau)=-2/\tau$ and
$g_z(\tau)=2\tau/(\tau^2+z^2)$.  The $h_0$ contribution converges uniformly in $r$ to a real constant $c_4$.  For $0<\sigma'\le B$ and $t^2\ge2(\tau_0^2+B^2)$,
\[
 |g_z(\tau_0\pm iy)|\le \frac{4(\tau_0+y)}{t^2+y^2}\le\frac{4(\tau_0+y)}{t^2}.
\]
Using the factor $2^{-my}$ from \eqref{eq:clock-halfplane-bound},
\[
 |R_m^{\pm,(z)}|\le\frac4{t^2}
 \left(\frac{\tau_0}{m\log2}+\frac1{(m\log2)^2}\right),
\]
so the corresponding series in \eqref{eq:Psi-r} is $O_B(t^{-2})$, uniformly in $r$.  Letting $r\uparrow1$ gives
\begin{equation}\label{eq:Psi-asymp}
 \Psi(z)=\Log(1+w(s))+c_4+O_B(t^{-2}).
\end{equation}

Combining \eqref{eq:transition-product}, \eqref{eq:transition-split},
\eqref{eq:smooth-difference}, and \eqref{eq:Psi-asymp} yields
\[
 \zeta_{\RVM}(s)
 =P_{\RVM}^{\crit}(1/2)e^{c_3+c_4}(1+w(s))(1+O_B(t^{-1})).
\]
This proves part (a) with
$Z_{\RVM}=P_{\RVM}^{\crit}(1/2)e^{c_3+c_4}>0$.

For part (b), define the holomorphic continuation of the Riemann--Siegel theta function on $|\Im u|<1/2$ by
\[
 \vartheta(u)=\frac1{2i}\left[\Log\Gamma\left(\frac14+\frac{iu}{2}\right)-
 \Log\Gamma\left(\frac14-\frac{iu}{2}\right)\right]-\frac u2\log\pi.
\]
The identity $\chi(1/2+iu)=e^{-2i\vartheta(u)}$ holds there by analytic continuation.  Stirling gives
$\mathcal N_{\RVM}(u)=\vartheta(u)/\pi+1+O(|u|^{-1})$ uniformly on closed substrips.  Since $s=1/2+ib$,
\[
 w(s)=\chi(s)(1+O(t^{-1})).
\]
On $a_0\le\sigma'\le b_0<1/2$, \eqref{eq:chi-size} makes $|\chi(s)|\le1/2$ for large $t$, and therefore
$1+w=(1+\chi)(1+O(t^{-1}))$.  Substitution into part (a) proves \eqref{eq:transition-chi}.  Conjugation gives $t<0$.
\end{proof}

\begin{corollary}[Vertical constant limit]\label{cor:vertical-limit}
For every fixed $\sigma>1/2$,
\begin{equation}\label{eq:vertical-limit}
 \zeta_{\RVM}(\sigma+it)\longrightarrow Z_{\RVM},\qquad
 \frac1{\zeta_{\RVM}(\sigma+it)}\longrightarrow L:=\frac1{Z_{\RVM}}
 \quad(|t|\to\infty).
\end{equation}
More precisely, with $\sigma'=\sigma-1/2$,
\begin{equation}\label{eq:w-decay}
 |w(\sigma+it)|=O_\sigma(|t|^{-\sigma'}),
\end{equation}
and hence
\begin{equation}\label{eq:FminusL-all}
 \frac1{\zeta_{\RVM}(\sigma+it)}-L
 =O_\sigma(|t|^{-\min(\sigma-1/2,1)}).
\end{equation}
\end{corollary}

\begin{proof}
Expanding the clock at $b=t-i\sigma'$ gives
$\Im\mathcal N_{\RVM}(b)=-(\sigma'/(2\pi))\log(t/(2\pi))+O_\sigma(t^{-1})$; hence
$|w|=\exp(2\pi\Im\mathcal N_{\RVM}(b))=O_\sigma(t^{-\sigma'})$.
Apply Theorem~\ref{thm:transition}(a) and invert for large $t$.
\end{proof}

\begin{theorem}[No ordinary absolutely convergent Dirichlet expansion]\label{thm:no-Dirichlet-RVM}
Neither $\zeta_{\RVM}$ nor $1/\zeta_{\RVM}$ is representable by an ordinary Dirichlet series that converges absolutely in any half-plane $\Re s>\sigma_0$.
\end{theorem}

\begin{proof}
Suppose first that $F(s)=1/\zeta_{\RVM}(s)=\sum_{n\ge1}a_n n^{-s}$ converges absolutely in some half-plane.  Choose a fixed
$\sigma>\max(\sigma_0,1/2)$.  Then
$F(\sigma+it)-L\to0$ by Corollary~\ref{cor:vertical-limit}, and absolute convergence permits termwise mean-square averaging:
\[
 \lim_{T\to\infty}\frac1{2T}\int_{-T}^T|F(\sigma+it)-L|^2dt
 =\sum_{n\ge1}|c_n|^2n^{-2\sigma},
\]
where $c_1=a_1-L$ and $c_n=a_n$ for $n\ge2$.  The left side is zero because the integrand tends to zero and is bounded outside a compact interval.  Hence every $c_n=0$, so $F\equiv L$ in that half-plane and therefore, by meromorphic continuation, identically.  This contradicts $F(1)=0$, since $\zeta_{\RVM}$ has a pole at $1$.

The same argument applied to $\zeta_{\RVM}(s)-Z_{\RVM}$ forces an assumed absolutely convergent Dirichlet series for $\zeta_{\RVM}$ to be the constant $Z_{\RVM}$, contradicting its pole and nontrivial zeros.
\end{proof}

We now return to the inverse transform.  The integral in \eqref{eq:Mcrit-inverse} is henceforth interpreted as the symmetric limit
$\lim_{U\to\infty}\int_{c-iU}^{c+iU}$.

\begin{theorem}[Unsmoothed synthetic Mertens bound]\label{thm:unsmoothed-RVM}
For every $x>1$ and every $c>1$, the symmetric Perron integral
\begin{equation}\label{eq:Mcrit-symmetric}
 M_{\RVM}^{\crit}(x)
 :=\lim_{U\to\infty}\frac1{2\pi i}
 \int_{c-iU}^{c+iU}\frac{x^s}{s\zeta_{\RVM}(s)}\,ds
\end{equation}
exists.  It is independent of $c>1/2$ when continued by the same symmetric formula, and for every $\varepsilon\in(0,1/2)$,
\begin{equation}\label{eq:Mcrit-unsmoothed-bound}
 \boxed{M_{\RVM}^{\crit}(x)=L+O_\varepsilon(x^{1/2+\varepsilon}).}
\end{equation}
In particular $M_{\RVM}^{\crit}(x)=O_\varepsilon(x^{1/2+\varepsilon})$.
\end{theorem}

\begin{proof}
Fix $c>1/2$.  By Corollary~\ref{cor:vertical-limit},
$F(c+it)-L=O_c(|t|^{-\delta_c})$ with $\delta_c=\min(c-1/2,1)>0$.  Therefore
\begin{equation}\label{eq:Mcrit-renorm}
 L+\frac1{2\pi i}\int_{(c)}(F(s)-L)\frac{x^s}{s}\,ds
\end{equation}
is absolutely convergent.  On the other hand the classical Bromwich kernel gives, for $x>1$,
\[
 \lim_{U\to\infty}\frac1{2\pi i}\int_{c-iU}^{c+iU}\frac{x^s}{s}\,ds=1.
\]
Splitting $F=L+(F-L)$ proves existence of \eqref{eq:Mcrit-symmetric} and equality with \eqref{eq:Mcrit-renorm}.

If $1/2<c<c'$, the integrand is holomorphic in the closed strip: the nontrivial poles of $F$ lie on $\Re s=1/2$, while $s=1$ is a zero of $F$.  On the horizontal sides $|F(s)-L|=O(1)$ by Theorem~\ref{thm:transition}(a), uniformly for $c\le\Re s\le c'$, and the constant part contributes $O(U^{-1})$; hence the horizontal integrals tend to zero.  Cauchy's theorem shows independence of $c$.

Take now $c=1/2+\varepsilon$.  From the absolutely convergent representation \eqref{eq:Mcrit-renorm},
\[
 |M_{\RVM}^{\crit}(x)-L|
 \le x^{1/2+\varepsilon}\frac1{2\pi}
 \int_{-\infty}^{\infty}\frac{|F(c+it)-L|}{|c+it|}\,dt,
\]
and the integral is finite by \eqref{eq:FminusL-all}.  This proves \eqref{eq:Mcrit-unsmoothed-bound}.
\end{proof}

\begin{remark}[Relation with the Gaussian theorem and with the Word note]
The Gaussian estimate of Proposition~\ref{prop:Gaussian} remains correct but is now weaker than Theorem~\ref{thm:unsmoothed-RVM}.  The latter is a theorem about the synthetic object only; it has no implication for the genuine Mertens function.  It also explains why the coefficient factorization proposed in Step~3.4 of the Word note cannot be the closure mechanism: $1/\zeta_{\RVM}$ has a well-defined symmetric Perron inverse but no absolutely convergent ordinary Dirichlet coefficient sequence.
\end{remark}

\section{Conclusions and theorem ledger}

The midpoint clock alone produces a coherent synthetic zeta system.  The following statements are proved without RH:
\begin{enumerate}[label=\textup{(\arabic*)}]
\item the Lambert midpoint ordinates and their canonical product are well defined;
\item $P_{\RVM}^{\crit}$ is entire and symmetric, and $\zeta_{\RVM}$ has precisely the prescribed simple critical zeros, a residue-one pole at $1$, and the usual trivial zeros;
\item $S_{\RVM}(N;u)=O_{a,b}(1+\log N)$ uniformly on compact positive frequency intervals, hence the point centroid and $M_H^{\RVM}(x)$ tend to zero;
\item the symmetric zero sum defining $J_{\RVM}$ exists compact-uniformly, and $J_{\RVM}^{*}$ and $\pi_{\RVM}$ are related by exact finite M\"obius inversion;
\item the nontrivial pole abscissa of $1/\zeta_{\RVM}$ is $1/2$ by construction;
\item the midpoint transition theorem gives $\zeta_{\RVM}=Z_{\RVM}(1+w)(1+O(1/|t|))$ on every fixed band to the right of $1/2$, and $Z_{\RVM}(1+\chi)(1+O(1/|t|))$ on closed substrips of $1/2<\Re s<1$;
\item consequently $1/\zeta_{\RVM}(\sigma+it)\to L=1/Z_{\RVM}$ for every fixed $\sigma>1/2$, neither $\zeta_{\RVM}$ nor its reciprocal is an absolutely convergent ordinary Dirichlet series in any right half-plane, and the symmetric Perron inverse exists with
$M_{\RVM}^{\crit}(x)=L+O_\varepsilon(x^{1/2+\varepsilon})$;
\item the exact quotient $\mathcal Q_{\RVM}=\zeta/\zeta_{\RVM}$ measures the transform-level defect from the classical Mertens system;
\item the continuous product transport from $\gamma_k^{\RVM}$ to $\gamma_k$ exists, has an explicit infinitesimal generator, total multiplier, endpoint expansion and residue transport;
\item $\zeta_{\linez}=\zeta$ is equivalent to RH.
\end{enumerate}

The following remain explicitly separate from that unconditional synthetic core: the square-root quantitative estimate for $J_{\RVM}$ and $\pi_{\RVM}$; unit-cell control and a positive $0/1$ integer realization; any boundedness or invertibility statement for the quotient $\mathcal Q_{\RVM}$ strong enough to transfer the synthetic bound to the genuine Mertens function; and any independent identification of the transported endpoint with the genuine arithmetic coefficients.  In particular, the three-step transport constructs the endpoint and proves the RH equivalence of its identification with $\zeta$, but it does not itself establish that identification.

\begin{thebibliography}{99}

\bibitem{FrancaLeClair2014}
G.~Fran\c{c}a and A.~LeClair,
\emph{Statistical and other properties of Riemann zeros based on an explicit equation for the $n$-th zero on the critical line},
arXiv:1307.8395v3 (2014).

\bibitem{DLMF-W}
NIST Digital Library of Mathematical Functions,
\emph{\S4.13 Lambert $W$-Function},
National Institute of Standards and Technology.

\bibitem{Titchmarsh1986}
E.~C. Titchmarsh,
\emph{The Theory of the Riemann Zeta-Function},
2nd ed., revised by D.~R. Heath-Brown,
Clarendon Press, Oxford, 1986.

\bibitem{Edwards1974}
H.~M. Edwards,
\emph{Riemann's Zeta Function},
Academic Press, 1974.

\bibitem{MontgomeryVaughan2007}
H.~L. Montgomery and R.~C. Vaughan,
\emph{Multiplicative Number Theory I: Classical Theory},
Cambridge Studies in Advanced Mathematics 97,
Cambridge University Press, 2007.

\bibitem{IwaniecKowalski2004}
H.~Iwaniec and E.~Kowalski,
\emph{Analytic Number Theory},
American Mathematical Society Colloquium Publications 53, 2004.

\bibitem{Hadamard1896}
J.~Hadamard,
\emph{Sur la distribution des z\'eros de la fonction $\zeta(s)$ et ses cons\'equences arithm\'etiques},
Bull. Soc. Math. France \textbf{24} (1896), 199--220.

\bibitem{BrouckeDebruyneVindas2020}
F.~Broucke, G.~Debruyne, and J.~Vindas,
\emph{Beurling integers with RH and large oscillation},
Adv. Math. \textbf{370} (2020), 107240.

\end{thebibliography}

\end{document}
